\subsection{System architecture}

\begin{figure}[H]
\centering
\begin{tikzpicture}[font=\scriptsize, >={Stealth[length=2mm]},
  box/.style={draw, rounded corners=1pt, align=center, minimum height=8mm, inner sep=3pt}]
  \node[box] (cam) {Camera\\(10~Hz)};
  \node[box, right=6mm of cam] (det) {ArUco\\detection};
  \node[box, right=6mm of det] (fus) {Spatial consensus +\\inverse-variance fusion};
  \node[box, right=6mm of fus] (kf) {Error-state\\Kalman filter};
  \node[box, below=10mm of kf] (guid) {Guidance\\(coarse standoff / fine align)};
  \node[box, left=14mm of guid] (ctl) {PBVS control +\\velocity feedforward};
  \node[box, left=14mm of ctl] (veh) {ArduSub\\+ vehicle};
  \node[box, below=8mm of ctl, dashed] (fsm) {Supervisory FSM: \textsc{idle} / \textsc{coarse} / \textsc{fine} / \textsc{docked}};
  \draw[->] (cam) -- (det);
  \draw[->] (det) -- (fus);
  \draw[->] (fus) -- (kf);
  \draw[->] (kf) -- node[right]{\tiny pose, velocity, health} (guid);
  \draw[->] (guid) -- node[above, yshift=1pt]{\tiny error} (ctl);
  \draw[->] (ctl) -- node[above, yshift=1pt]{\tiny cmd\_vel} (veh);
  \draw[dashed, ->] (veh.north) -- node[left]{\tiny camera moves with vehicle} (cam.south);
  \draw[dashed, ->] (fsm.north) -- (ctl.south);
  \draw[dashed, ->] (fsm.east) -| node[pos=0.35, below]{\tiny engage / handoff / demote} (guid.south);
\end{tikzpicture}
\caption{System architecture. Perception runs entirely in the loop; the
supervisory state machine sequences the mission and gates the controllers
on the filter's health signal. The feedback path from vehicle to camera is
the ego-motion coupling analysed in Section~\ref{sec:results-a}.}
\label{fig:arch}
\end{figure}

The system is a perception, guidance and control pipeline for the terminal stage of docking. Perception estimates the dock pose from passive ArUco markers and publishes a filtered pose, a dock velocity estimate, and a semantic health signal. Guidance converts the dock pose and the vehicle pose into a body-frame error to a target point. Control regulates that error with position-based visual servoing (PBVS) and writes a body velocity command. A supervisory state machine sequences the mission through \textsc{idle}, \textsc{coarse} (drive to a standoff point on the dock entry axis), \textsc{fine} (align and advance into the entry) and \textsc{docked}, with demotion paths on loss of the perception signal.

PBVS was selected over image-based servoing because the marker pipeline already produces a full metric 6-DOF pose: fusing multiple markers of known geometry yields a calibrated dock pose, so servoing in pose space costs nothing extra and keeps the guidance geometry (standoff points, entry axis, capture tolerances) expressible in physical units \cite{chaumette2006visual}.

The scope boundary is the optical terminal stage: trials begin with the dock inside the camera field of view at several metres range. Long-range acquisition (acoustic homing or an active light cue) is assumed to have delivered the vehicle to this point and is out of scope.

\begin{figure}[H]
\centering
\begin{tikzpicture}[font=\scriptsize, >={Stealth[length=2mm]}, node distance=16mm,
  st/.style={draw, rounded corners=3pt, align=center, minimum height=9mm, inner sep=4pt}]
  \node[st] (idle) {\textsc{idle}\\ \tiny re-arm, hold};
  \node[st, right=of idle] (coarse) {\textsc{coarse}\\ \tiny drive to standoff};
  \node[st, right=of coarse] (fine) {\textsc{fine}\\ \tiny align, advance};
  \node[st, right=of fine] (dock) {\textsc{docked}\\ \tiny disarm};
  \draw[->] (idle) -- node[above, align=center]{\tiny engage} (coarse);
  \draw[->] (coarse) -- node[above, align=center]{\tiny handoff gate:\\ \tiny pos + yaw tol,\\ \tiny velocity-matched,\\ \tiny debounced} (fine);
  \draw[->] (fine) -- node[above, align=center]{\tiny seated $\leq$ 0.10 m,\\ \tiny debounced} (dock);
  \draw[->] (fine.south) to[bend left=25] node[below, align=center]{\tiny demote: loss $>$ 2 s or range $>$ 1.5 m} (coarse.south);
  \draw[->] (coarse.north) to[bend right=30] node[above]{\tiny disengage} (idle.north);
  \draw[->] (dock.north) to[bend right=45] node[above]{\tiny disengage (re-arm on \textsc{idle} entry)} (idle.north);
\end{tikzpicture}
\caption{Docking state machine with the transition gates annotated.
Demotion paths make perception loss recoverable rather than terminal.}
\label{fig:fsm}
\end{figure}
